\section{Results and discussion}

\subsection{Closed-spectrum surrogate modeling}

The closed-spectrum benchmark establishes that the direct tight-binding workflow
is physically calibrated at the level needed for a conservative surrogate
study, but it does not establish an ML novelty claim. Two scale checks play
different roles. The circular Bessel-scale check tests the low-energy convention
and circular calibration; it gives an error of about $2.03\%$. The fixed-shape
scaling diagnostic tests whether the expected leading $a^{-2}$ confinement
scaling is consistent across the five tested sizes; it finds a maximum
deviation of about $2.12\%$ for $(E_0+4)a^2$. These are different diagnostics
and should not be interpreted as independent measurements of the same error.
Together they support the benchmark convention and leading confinement scale,
not a nontrivial shape-design rule.

Within this dataset, the physics-informed Ridge model remains the preferred
surrogate. A small multilayer perceptron with physics-motivated features passes
the cell-level success rule in $10$ of $16$ main comparison cells, but only in
$3$ of $8$ leave-one-aspect-out cells and $7$ of $8$ leave-one-aspect-ratio-out
cells. Thus the preregistered robust-success criterion is not met. The numerical
scale of the Ridge error is already small: the relative MAE is
$2.01$--$2.60\%$ of the characteristic $E_{\mathrm{kin}}$ for $E_0$ and
$0.78$--$1.15\%$ of the characteristic $dE_1$ gap. The MLP has individual
winning cells, including a large favorable improvement for one $E_0$ protocol,
but it also has strongly negative cells for $dE_1$. The safe conclusion is that
the MLP does not robustly outperform a physically motivated linear baseline in
the structured validation used here.

\begin{figure*}[t]
    \centering
    \includegraphics[width=0.88\textwidth]{surrogate_validation_summary.pdf}
    \caption{
    Surrogate-model diagnostic from existing closed-spectrum outputs. The left
    panel summarizes the MLP success counts under progressively more structured
    validation splits. The right panel summarizes the Ridge relative-MAE ranges
    used as the conservative baseline scale. The figure is a visualization of
    repository outputs only; it does not introduce new training runs.
    }
    \label{fig:surrogate-diagnostic}
\end{figure*}

This result matters for later inverse-screening attempts. If the surrogate is
not a robust model result, then surrogate-generated candidates must be treated
only as proposals. Direct Kwant verification and comparison with physics
baselines remain the source of any claim.

\subsection{Spectral objectives and inverse-screening limitations}

The first inverse-screening direction used the ratio
$Q=dE_1/(E_0+4)$ as a one-shot surrogate-guided objective inside the verified
training-domain ranges. The report explicitly treats Ridge predictions as
candidate generators only, with final candidate and baseline values coming from
direct Kwant calculations or already computed training rows. The outcome is
negative for each fixed $n$ class. The symmetry-optimum follow-up further shows
that the earlier top inverse-screening candidate and the isotropic same-$n$
baseline are the same discrete geometry for every tested $n$.

The preregistered S-objective line tests
$S=(E_2-E_1)/(E_0+4)$ under a frozen rule set. Its own report emphasizes that
observed doublet splitting is not a discovery: the question is whether
surrogate-guided screening beats strong baselines. It does not. At the primary
setting $\alpha=0.95$, $0$ of $4$ fixed $n$ values pass; at the secondary
setting $\alpha=0.90$, the result also remains $0$ of $4$. The strongest
baseline is the simple anisotropy heuristic for the reported rows.

Together these results limit a broad class of spectral ratio objectives. Gap
ratios and doublet-splitting measures can encode obvious geometric trends such
as isotropy, aspect ratio, and symmetry lifting. Without a stronger baseline
ladder, they should not be used as evidence for inverse design or new
shape-dependent physics.

\subsection{Magnetic-response ranking limitation}

The magnetic-response ranking analysis was a direct-Kwant check, not an ML task
and not an extension of the Q/S objectives. It tested four representative
geometries at sizes $a=\{30,36\}$, with weak-field flux parameters
$\alpha=(0.0,0.00125,0.0025,0.005)$ and diagnostic values
$\alpha=(0.01,0.02,0.04)$. Numerical checks passed: the $\alpha=0$ spectra
were reproduced, gauge invariance was verified with maximum reported
difference $1.865174681370263\times10^{-14}$, Hermiticity and finite sorted
spectra checks passed, and eigenvalue imaginary parts were controlled.

The proposed geometry-ranking signal is not supported. The analysis found $0$
thresholded pairwise ranking crossovers and $0$ crossovers surviving the
size-stability filters. Three robustness-divergence candidates were identified,
but $0$ survived those filters. The apparent candidates are attributable to
circle or ellipse symmetry-lifting baselines. The field-dependent calculation is
numerically controlled, but this tested protocol does not produce a non-baseline
geometry-ranking signal.

\subsection{Continuum-reference limitation}

The continuum-residual analysis asked whether tight-binding spectra contain a
nontrivial superellipse-exponent-dependent finite-lattice convergence residual
after subtracting ordinary continuum, Weyl, perimeter, pixelation, and
discretization baselines. The protocol deliberately forbids weak claims based
on the leading $1/a^2$ scale, generic boundary effects, ML accuracy, or
largest-size tight-binding extrapolation. Its first check is a circular Bessel
anchor for $n=2.0$, $r_{\mathrm{AR}}=1.0$, with sizes
$a=\{24,30,36,48,60,72,96\}$ and levels $E_0$ through $E_5$.

That anchor passes its stability review, which is useful but narrow. It
calibrates the low-energy circular case; it does not establish a
shape-sensitive continuum residual for $n\neq2.0$. The finite-difference
reference step then validates an embedded-mask Dirichlet Laplacian on grids up
to $N_{\mathrm{grid}}=251$, where the selected circular ground-state relative
error versus the Bessel value is $0.0052206740942014095$. This is adequate as
an initial reference check, but the later $N_{\mathrm{grid}}=501$ convergence
check changes the conclusion for shape-sensitive use.

The N501 check uses the near-geometric refinement triple $(126,251,501)$,
corresponding approximately to $h=\{0.016,0.008,0.004\}$. Although the raw
circle ground-state relative error versus Bessel improves to
$0.0026734921480308575$, the non-circular reference uncertainties remain too
large for the downstream residual claim:
$0.0047216251159203785$ for the $n=1.2$ level-zero reference and
$0.0021923716263753637$ for the $n=4.0$ level-zero reference. Both exceed the
$0.001$ reference-uncertainty criterion. The safe conclusion is not that FD is
generally invalid, but that this embedded-mask FD reference should not be
treated as continuum truth for the intended shape-sensitive residual.

\subsection{TB-only finite-size scaling limitation}

After the FD reference was found insufficient for the continuum-residual
analysis, the TB-only self-scaling diagnostic used direct tight-binding
ground-state values and geometric diagnostics only. No FD reference values were
loaded, no TB-vs-FD residuals were computed, and no ML, Q, or S objectives were
used. The domain was $r_{\mathrm{AR}}=1.0$, $n=(1.2,2.0,4.0)$, and
$a=(24.0,30.0,36.0,48.0,60.0,72.0,96.0)$.

The circular calibration is good: for $n=2.0$, the effective-radius model gives
a relative error of $7.21\times10^{-4}$ against
$j_{01}^2=5.783185962946783$. Fitted effective-radius parameters are also
reported for the three $n$ values. The limiting coefficients are approximately
$8.03$, $5.779$, and $5.049$, with offsets approximately $0.279$, $0.277$, and
$0.130$ for $n=1.2,2.0,4.0$, respectively. These values show ordinary shape
dependence of continuum-like eigenvalues and finite-size offsets. They are not,
by themselves, a supported shape-residual signal.

The potential residual exponent is unstable. The reported power-law values are
approximately $p_{\mathrm{TB}}=0.10,0.53,1.05$ for $n=1.2,2.0,4.0$, but the
power-law leave-one-size-out stability test is not accepted for any of the
three $n$ values. The corresponding leave-one-size-out $p$ ranges are about
$0.67$, $1.73$, and $1.01$, respectively, and the $n=1.2$ fit sits at the lower
bound of the implemented power-law fit. The existing diagnostic therefore
reports $p_{\mathrm{TB}}$ as an instability indicator, not as a physical
exponent with a defensible uncertainty interval. Simple geometric baselines do
not cleanly explain a robust residual, but the residual itself is not stable
enough to support a shape-residual claim. The supported scientific result is
inconclusive.

\begin{figure*}[t]
    \centering
    \includegraphics[width=0.90\textwidth]{continuum_tb_diagnostics.pdf}
    \caption{
    Continuum-reference and TB-only scaling diagnostics from existing outputs.
    The FD reference uncertainties for non-circular shapes exceed the
    $10^{-3}$ criterion used before shape-sensitive residual analysis. The
    TB-only power-law exponents are accompanied by large leave-one-size-out
    ranges and are not treated as physical exponents.
    }
    \label{fig:continuum-tb-diagnostic}
\end{figure*}

\subsection{Open-transport limitation}

The open-transport diagnostic changes the observable from closed spectra to
two-terminal conductance. This is physically richer, but it introduces a new
baseline: contact geometry. The tested setup uses square-lattice strip leads,
the spinless convention $G=T$, a superellipse dot with $a=30.0$ and
$r_{\mathrm{AR}}=1.0$, lead widths $W=4$ and $W=10$, and an energy window
$[-3.8,-3.0]$ sampled at $241$ points. The curve distance is the raw normalized
RMS distance defined in Eq.~(\ref{eq:transport-distance}); all raw distances are
computed on the same energy grid.

The central result is that contact width dominates the tested shape effect. For
$n=2.0$, the raw distance between $W=4$ and $W=10$ conductance curves is
$0.6589809313045025$. The mean pairwise shape distances are
$0.5810925658201803$ at $W=4$ and $0.5122909816541278$ at $W=10$. The
shift/rescale diagnostic gives $0.24323761388169962$ for the contact comparison
and $0.49646370215309615$ and $0.2802443376667481$ for the two shape
comparisons. Mode thresholds do not explain the features, the straight-channel
baseline is computed but does not remove the result, and an energy-shift/rescale
collapse does not remove it either. The limitation is therefore not a trivial
numerical failure. It is a physical interpretation problem: in this open-dot
geometry, lead-dot coupling is at least as important as the shape parameter
being tested.

\begin{figure*}[t]
    \centering
    \includegraphics[width=0.70\textwidth]{transport_effect_summary.pdf}
    \caption{
    Compact summary of the open-transport diagnostic. In the tested setup, the
    contact-width effect for $n=2.0$ is larger than the raw normalized
    shape-effect summaries at both lead widths. The plotted values are rounded
    from the source inventory and use Eq.~(\ref{eq:transport-distance}).
    }
    \label{fig:transport-effects}
\end{figure*}

\begin{table*}[t]
\caption{
Merged summary of tested analyses, diagnostic values, limiting baselines, and
safe conclusions. Exact values and provenance are listed in the source
inventory; this table uses rounded values only to keep the audit readable.
}
\label{tab:hypothesis-summary}
\scriptsize
\setlength{\tabcolsep}{2pt}
\begin{tabular}{@{}llllll@{}}
\toprule
\tcell{0.12\textwidth}{Analysis} &
\tcell{0.13\textwidth}{Observable / protocol} &
\tcell{0.17\textwidth}{Key diagnostic value} &
\tcell{0.17\textwidth}{Main baseline or limitation} &
\tcell{0.11\textwidth}{Result} &
\tcell{0.22\textwidth}{Safe conclusion} \\
\midrule
\tcell{0.12\textwidth}{Closed-spectrum Ridge/MLP} &
\tcell{0.13\textwidth}{$E_0$, $dE_1$ under structured validation} &
\tcell{0.17\textwidth}{Ridge relative MAE $2.01$--$2.60\%$ for $E_0$ and $0.78$--$1.15\%$ for $dE_1$; MLP success $10/16$, $3/8$, $7/8$} &
\tcell{0.17\textwidth}{Physics-informed Ridge and split-aware validation} &
\tcell{0.11\textwidth}{ML advantage not established} &
\tcell{0.22\textwidth}{Ridge remains the preferred conservative surrogate; direct Kwant remains the reference.} \\
\tcell{0.12\textwidth}{Q inverse screening} &
\tcell{0.13\textwidth}{$dE_1/(E_0+4)$} &
\tcell{0.17\textwidth}{Success false for each fixed $n$; top candidate equals same-$n$ isotropic baseline} &
\tcell{0.17\textwidth}{Isotropy and same-$n$ baselines} &
\tcell{0.11\textwidth}{Not supported} &
\tcell{0.22\textwidth}{Q candidates do not justify inverse-design claims.} \\
\tcell{0.12\textwidth}{S objective} &
\tcell{0.13\textwidth}{$(E_2-E_1)/(E_0+4)$} &
\tcell{0.17\textwidth}{$0/4$ pass at $\alpha=0.95$ and $0/4$ at $\alpha=0.90$} &
\tcell{0.17\textwidth}{Simple anisotropy and strongest feasible baselines} &
\tcell{0.11\textwidth}{Not supported} &
\tcell{0.22\textwidth}{S behaves as an anisotropy diagnostic in the tested domain.} \\
\tcell{0.12\textwidth}{Magnetic-response ranking} &
\tcell{0.13\textwidth}{Weak-field level rankings} &
\tcell{0.17\textwidth}{Thresholded crossovers $0$; surviving crossovers $0$; robustness divergences surviving filters $0$} &
\tcell{0.17\textwidth}{Symmetry lifting and zero-field/aspect-ratio ranking} &
\tcell{0.11\textwidth}{Baseline limited} &
\tcell{0.22\textwidth}{Field-dependent spectra are controlled numerically, but no non-baseline design signal is supported.} \\
\tcell{0.12\textwidth}{FD continuum reference} &
\tcell{0.13\textwidth}{Embedded-mask FD reference for shape-sensitive residuals} &
\tcell{0.17\textwidth}{Non-circular level-zero relative uncertainties $0.00472$ and $0.00219$ exceed $0.001$} &
\tcell{0.17\textwidth}{Reference-solver uncertainty} &
\tcell{0.11\textwidth}{Inconclusive} &
\tcell{0.22\textwidth}{Current FD reference should not be treated as continuum truth for shape-sensitive residuals.} \\
\tcell{0.12\textwidth}{TB-only scaling} &
\tcell{0.13\textwidth}{Ground-state finite-size scaling} &
\tcell{0.17\textwidth}{$n=2$ calibration $7.21\times10^{-4}$; $p_{\mathrm{TB}}=0.10,0.53,1.05$ with unstable LOO ranges} &
\tcell{0.17\textwidth}{Unstable residual exponent and insufficient surviving signal} &
\tcell{0.11\textwidth}{Inconclusive} &
\tcell{0.22\textwidth}{The diagnostic does not support a continuum-residual claim.} \\
\tcell{0.12\textwidth}{Open transport diagnostic} &
\tcell{0.13\textwidth}{$G(E)$ curves, $W=4,10$} &
\tcell{0.17\textwidth}{Contact distance $0.659$; shape distances $0.581$ and $0.512$} &
\tcell{0.17\textwidth}{Contact-width dependence and lead-dot coupling} &
\tcell{0.11\textwidth}{Contact dominated} &
\tcell{0.22\textwidth}{Transport interpretation requires tighter lead-dot coupling control before a shape claim.} \\
\bottomrule
\end{tabular}
\end{table*}

\clearpage
